\chapter{What can be proved, and under which assumptions}
\label{ch:proof}

\section{Giving a correctness claim a precise target}
\label{sec:proof-target}
The second-circular exercise leaves us with two questions that sound similar but
have different evidentiary requirements. Does the program compute the proposed
gift rule correctly? Does the proposed gift rule express a justified reading of
the circular? The first question becomes mathematical once the rule and its
inputs have precise meanings. The second includes deciding which meanings the
words, context and authority support. A proof of the first cannot silently
include a proof of the second.

To see why, suppose a programmer writes a formula that checks gifts promoting a
specific product and omits the product-type route. A proof assistant can verify
that an implementation returns exactly that formula's value for every Boolean
assignment. The proof is correct about its target. The target is incomplete
relative to the proposed source interpretation. Making the proof longer, using a
more powerful solver or verifying another implementation will not add the
missing route. The source-to-formula mapping has to be challenged directly.

This is not a special weakness of legal language. Formal verification always
depends on the proposition that was formalized. The virtue of a formal method
is that it makes the conditional claim exact and inspectable. It does not remove
the need to establish that the proposition is the one the institution intended
to rely on. Catala's separation between legal authoring and compiler semantics
is useful precisely because it keeps these responsibilities distinct
\citep[core language and compilation sections]{catala2021}.

Let $E$ denote a versioned source packet and $S$ a controlled specification
adopted for a declared interpretation. Let $F$ be a RuleIR program and $J$ the
generated Java artifact. Let $A$ contain the explicit vocabulary definitions,
domain restrictions and operational assumptions. The input domain $D_A$ consists
of facts and histories allowed by those assumptions. We can now state a proposed
preservation obligation:
\begin{equation}
 \forall x\in D_A,\qquad
 \pi\bigl(\eval_{\CNL}(S,x)\bigr)
 =\pi\bigl(\eval_{\IR}(F,x)\bigr)
 =\pi\bigl(\eval_{J}(J,x)\bigr).
 \label{eq:preservation-chain}
\end{equation}
The projection $\pi$ selects the semantic result components that the contract
requires to agree. These include the decision category and the specified
explanation and evidence behavior. They exclude fields whose purpose is to
identify different backends. Chapter~\ref{ch:verification} makes that projection
concrete. Equality in equation~\eqref{eq:preservation-chain} is a proposed formal
target for a supported language fragment, not a theorem already proved about
the delivered application.

The claim that $S$ is justified by $E$ remains a separate interpretation claim.
It can be supported by source anchors, explicit alternatives, reviewed
definitions, adversarial cases and an attributable judgment. When some part of
$E$ has itself been translated into a formal language, a theorem can relate
those formal objects. The choice that made the formal source representation
faithful still needs justification. Repeating that translation boundary does not
eliminate it.

\section{Quantifiers change which entity carries the obligation}

The phrase ``for each'' introduces a choice that a Boolean expression can easily
hide. To see it without assuming knowledge of predicate logic, consider a
synthetic requirement: every relevant benefit in an offer must satisfy a specified
condition. Let $B(o)$ be the finite set of benefit components in offer $o$, and
let $P(b)$ mean that component $b$ has the property under investigation. The
universal reading is
\begin{equation}
 \forall b\in B(o),\ P(b).
 \label{eq:universal-benefit}
\end{equation}
It requires every component to satisfy the condition. An existential reading,
\begin{equation}
 \exists b\in B(o),\ P(b),
 \label{eq:existential-benefit}
\end{equation}
requires only one. These expressions are not two styles of writing the same rule.
For an offer with two components, one satisfying $P$ and one not, the universal
reading is false and the existential reading true.

The distinction becomes more consequential when an exception is involved. Suppose
$G(b)$ identifies a gift and $D(b)$ an excepted discount. A component-level
prohibition trigger might be
\begin{equation}
 \exists b\in B(o),\ G(b)\land\neg D(b).
 \label{eq:component-exception}
\end{equation}
An alternative whole-offer encoding might be
\begin{equation}
 \left(\exists b\in B(o),\ G(b)\right)
 \land\neg\left(\exists b\in B(o),\ D(b)\right).
 \label{eq:offer-exception}
\end{equation}
The first asks whether any component is a non-excepted gift. The second suppresses
the trigger whenever any component is a discount. A voucher combined with a fee
waiver distinguishes them immediately: the second formula can let the waiver
cancel the voucher's effect, while the first continues to assess the voucher.
These are proposed logical readings for analysis, not a legal determination about
mixed offers under 23EC46.

A source interpretation must decide what the variable ranges over before a proof
about either formula is useful. Is the subject a benefit, a contract, a campaign,
a transaction or a client relationship? Are components defined by economic
substance, contractual terms or an operational product record? The formal
notation exposes this question because $B(o)$ cannot remain undefined in an
implementation-ready specification.

The current RuleIR fragment has scalar facts rather than unrestricted quantified
collections. A reviewed adapter can create one snapshot per component and a host
orchestration can aggregate results under an explicit rule. That aggregation is
another specification obligation. It must define how true, false, unknown,
conflict and out-of-scope component results combine, and how the host knows the
component list is complete.

An empty list deserves attention. In classical logic, a universal statement over
an empty set is true and an existential statement is false. That convention is
mathematically coherent, but a missing component inventory must not be treated as
an established empty offer. The adapter needs an evidence status for collection
completeness. Otherwise a data omission can pass a universal check by vacuity.
The same pattern occurs when no argument extensions are returned or no test cases
are selected: an empty computation can look successful unless existence and
coverage are checked separately.

Unknown classification also affects aggregation. If one component is known to
trigger a prohibition and another is unknown, the selected logical aggregation
may already establish a trigger. Yet the unknown component may remain relevant
to another obligation or to the completeness of the offer review. A top-level
known result should not erase the record of missing facts. The distinction between
missing and blocking inputs in the runtime supports this more informative report.

\section{Necessary conditions, sufficient conditions and permissions}

Another frequent translation error changes the direction of implication. Consider
a synthetic statement that a process may use a streamlined route only if consent
has been obtained. Let $S$ mean use of the streamlined route and $C$ mean established
consent. The necessary-condition reading is $S\Rightarrow C$: using the route
requires consent. It does not follow that $C\Rightarrow S$ or that consent alone
authorizes the route. Other eligibility conditions may still apply.

A truth-table example makes the difference concrete. With consent true and an
unmet financial condition, the necessary-condition statement is not violated by
refusing the streamlined route. An implementation that treats consent as sufficient
would proceed. Both programs may mention the same two words, but their decisions
represent different logical relationships.

Likewise, an exception to one prohibition is not a general permission. If the
selected condition $R$ is false, we know only that this particular trigger is not
satisfied under the accepted facts and semantics. We do not infer that every
other prohibition is false or every positive obligation fulfilled. The host must
compose independently applicable controls without upgrading the absence of one
trigger into universal clearance.

In a closed, complete formal theory, one can define permission in terms of the
absence of a prohibition if that is the intended logic. Legal systems and practical
source packets often do not provide the completeness needed for that move.
Explicit permission, lack of evidence of prohibition and a bank's operational
choice should therefore have different representations. The normative-language
literature examined in Chapter~\ref{ch:languages} provides several treatments;
RuleIR's narrow condition results do not choose a universal deontic logic.

The words ``unless'' and ``except'' also require a structural decision. A simple
propositional rewriting may express an exception as a negated condition, while a
default system represents a defeasible conclusion that can be overridden. These
can agree on a selected complete fact pattern and differ under missing evidence,
conflicting exceptions or additional rules. The translator must identify which
behavior the approved specification requires.

For the circular example, using $\neg D$ inside the selected body is an explicit
modeling decision tied to the reviewed factual definition of discount. If $D$ is
unknown, the program should not act as though the exception is certainly absent.
The three-valued semantics makes that uncertainty visible. A closed-world database
query that treats a missing discount row as false would encode a different
assumption, which requires separate evidence and approval.

\section{A proof obligation should expose its assumptions}

A useful proof statement includes the source-independent mathematical property,
the accepted semantic definitions and the domain assumptions. For example:
for every well-typed Boolean expression in a specified grammar and every normalized
fact assignment without conflict, the reference and target evaluators produce
the same three-valued result. This statement is precise enough to prove by
structural induction. It does not include the claim that a particular English
paragraph was mapped to the right expression.

The interpretation assumption can be represented separately as a reviewed
proposition: under source packet $S$, profile $P$ and definitions $D$, expression
$e$ represents the selected requirement. A preservation proof then establishes
that target program $J(e)$ behaves like $e$. Composing these steps is useful, but
the first premise remains a source and review judgment unless a restricted
controlled language and its translation semantics are themselves the authoritative
input.

A proof assistant will faithfully carry an unjustified premise into a theorem.
If we assume that every cashback is a discount, it may prove the corresponding
exception applies. The theorem is conditional on that assumption. It is not a
discovery that the assumption is legally correct. The proof report must display
such premises in language the reviewer can understand.

The same principle applies to abstraction. A model may replace a complex offer
with six Boolean inputs. The proof can establish behavior for every assignment
of those inputs. To apply it to the real offer, the adapter must establish that
the six inputs faithfully represent the relevant facts and that the omitted
features cannot change the selected requirement within the declared scope.
That relationship is an abstraction argument, not an automatic consequence of
exhaustive truth-table testing.

The workbench should therefore attach an assumption register to every formal
result. Each assumption names its provenance, operational meaning and review
status. A proof that depends on an unresolved definition is retained as a valid
conditional engineering result, while release remains blocked. This preserves
the value of formal work without allowing its precision to conceal an uncertain
starting point.

\section{Expressions, types and meanings}
\label{sec:proof-language}
A formal language is a collection of permitted expressions with rules for
deciding what they mean. It need not look like a conventional programming
language. A decision table, a restricted English sentence or a state-transition
diagram can have formal semantics if each permitted construct has a precise
interpretation. An unrestricted paragraph usually has more interpretive freedom
than such a language allows.

The simplest fragment needed for the gift example contains Boolean literals,
references to named facts, conjunction, disjunction and negation. In a compact
grammar, an expression $e$ is generated by
\begin{equation}
 e ::= \mathsf{true}\mid\mathsf{false}\mid\mathsf{fact}(a)
       \mid\mathsf{not}(e)\mid\mathsf{all}(e_1,\ldots,e_k)
       \mid\mathsf{any}(e_1,\ldots,e_k),\quad k\geq 1.
 \label{eq:boolean-grammar}
\end{equation}
Here $a$ is the identifier of a declared Boolean fact. The grammar specifies
which trees are valid expressions; it does not decide whether the fact denotes
a legally appropriate classification. That meaning belongs to the declaration
and its source-linked interpretation.

Types prevent several avoidable errors. A Boolean condition is different from
a monetary amount. A date is different from a number of days. Money denominated
in HK cents is different from a foreign-currency amount before conversion.
The current RuleIR scalar types are Boolean, integer, HKD money and date. A
comparison accepts compatible ordered values. Addition preserves a numeric
type. No implicit conversion turns a missing Boolean into zero or a dollar
amount into cents.

Typing is valuable because it can reject an entire class of malformed
specifications before any client case is evaluated. It is also limited. Two
variables can both have the Boolean type while referring to different subjects
or definitions. A type checker cannot infer that an offer-level discount flag
was wrongly substituted for a component-level classification unless the type
and declaration system explicitly represent that distinction. The proposed
extension should therefore retain subject, scope and definition identities
alongside the scalar type.

An expression's semantics is its evaluation rule. For known Boolean inputs,
\texttt{all} is true exactly when every operand is true, and \texttt{any} is true
when at least one operand is true. The gift body is a tree whose root is an
\texttt{all}, with children for the gift fact, the two-route disjunction and
the negated exception. The tree makes parentheses explicit. A source sentence
can be attached to the root and narrower evidence anchors to its children.

The chosen language is deliberately finite and acyclic. A reference cycle such
as rule A depending on B while B depends on A is rejected in the current
profile. This avoids an undefined recursive evaluation. Richer legal formalisms
can assign meaning to recursive or mutually defeating rules, but doing so
requires a specified semantics. A compiler must not guess that an unsupported
recursive construct means the same thing as the current acyclic profile.

\section{Unknown is an information state}
\label{sec:proof-unknown}
The bank often lacks a decisive fact when a rule is evaluated. In the gift
example, the evidence may not establish whether an additional return is a gift.
The program should distinguish that absence from evidence that no gift exists.
We therefore need an unknown value, written $\Unknown$, alongside known true
$\True$ and known false $\False$.

The current Boolean operations follow strong-Kleene tables. The name identifies
a precise three-valued calculation, which can be taught through the two cases
already encountered. If a product-specific promotional link is true, then the
disjunction of that link with an unknown type link is true. If the specific
link is false and the type link is unknown, the disjunction is unknown. A true
route suffices; an absent route cannot settle the other one.

\begin{table}[htbp]
\centering
\begin{tabular}{cc|cc}
\toprule
$a$ & $b$ & $a\land b$ & $a\lor b$\\
\midrule
$\True$ & $\True$ & $\True$ & $\True$\\
$\True$ & $\False$ & $\False$ & $\True$\\
$\True$ & $\Unknown$ & $\Unknown$ & $\True$\\
$\False$ & $\False$ & $\False$ & $\False$\\
$\False$ & $\Unknown$ & $\False$ & $\Unknown$\\
$\Unknown$ & $\Unknown$ & $\Unknown$ & $\Unknown$\\
\bottomrule
\end{tabular}
\caption{The omitted rows follow by exchanging the operands. Negation swaps
true and false and leaves unknown unchanged. The table describes logical values;
evidence conflict and errors are handled separately.}
\label{tab:kleene}
\end{table}

An established exception can make the gift conjunction false even when another
body fact is unknown, because one false conjunct is enough. That is the logic
behind g16. It is not a general instruction to ignore missing information.
Unknown scope is handled before the body, and conflicting evidence has its own
precedence. The evaluation result must retain which question was answered and
why the remaining unknown body facts did not change that answer.

A convenient symbolic encoding represents a three-valued Boolean by two ordinary
Boolean indicators: one means ``known true'' and the other ``known false.'' The
unknown value has neither indicator. If $t_a$ and $f_a$ encode $a$, then
\begin{align}
 t_{a\land b}&=t_a\land t_b,& f_{a\land b}&=f_a\lor f_b,\label{eq:kleene-and}\\
 t_{a\lor b}&=t_a\lor t_b,& f_{a\lor b}&=f_a\land f_b,\label{eq:kleene-or}\\
 t_{\neg a}&=f_a,& f_{\neg a}&=t_a.\label{eq:kleene-not}
\end{align}
The domain excludes $(t_a,f_a)=(1,1)$. That pair would assert both truth and
falsity, which belongs to a different conflict mechanism in the current design.
The equations reproduce Table~\ref{tab:kleene}: for example, true and unknown
have true indicators 1 and 0, so their conjunction cannot be known true; neither
has a false indicator, so the result is unknown.

This encoding allows a solver to search disagreements while preserving missing
information. Replacing each unknown by an arbitrary ordinary Boolean inside
one execution would ask a different question. It could produce a definite answer
chosen from a possible completion without showing that other completions differ.
A comparison tool must name the semantics it encoded, not merely report that it
used a solver.

There is a subtler distinction between compositional three-valued evaluation
and evaluating every Boolean completion. Let $p$ be unknown. Strong-Kleene
evaluation gives
\begin{equation}
 p\lor\neg p=\Unknown\lor\Unknown=\Unknown.
 \label{eq:unknown-tautology}
\end{equation}
Under every consistent Boolean completion, however, the expression is true:
if $p$ is true the first term holds, and if $p$ is false the second holds. A
completion-based semantics can therefore return true. Neither calculation is
ambiguous once its rules are specified; they are different semantics.

For the read-once gift body, the retained truth oracle's completion method agrees
with the compositional calculation. That particular equivalence does not make
the completion oracle a general replacement for RuleIR semantics. A future
ensemble should treat an unexpected disagreement as a possible semantics
mismatch before assuming that one interpreter contains a bug. This example is
also a reason to test repeated variables when extending the comparison fragment.

\section{Conflict, scope and errors are not extra Boolean values}
\label{sec:proof-tags}
Unknown truth and conflicting evidence should not be combined casually into one
truth table. The current runtime first normalizes records by validity,
recording time and explicit supersession. Two admissible non-superseded records
with incompatible values produce conflict. That result records a problem in
the evidence supporting a fact. It does not ask the Boolean evaluator to choose
which record is right.

The static dependency closure contains every fact reachable from the requested
rule through its scope, body, referenced rules and exception branches. Conflict
in that closure prevents evaluation from treating a convenient branch as a way
to evade contradictory evidence. A conflict in an unrelated declared fact does
not automatically block the requested rule. This distinction makes the
conservative policy bounded and reproducible.

After successful structural validation and conflict handling, scope determines
whether the body is an appropriate question. A known-false scope produces out
of scope. An unknown scope produces unknown with a scope reason. Only a
known-true scope permits the in-scope body result. Thus an out-of-scope result
is a statement about the relation between the query and the selected rule,
not a false answer to the rule's substantive condition.

Malformed data and unsupported constructs introduce a further distinction.
An input with an invalid date or duplicated conflicting identifier may be an
error, not legal uncertainty. If the system cannot parse a request, it cannot
honestly report that the regulatory condition is unknown on a well-defined
request. The API should identify the boundary error and the location of the
malformation. Operationally both may prevent automatic action, but their causes
and repair routes differ.

We can therefore describe the result as a tagged value whose tag identifies
the kind of answer. The current contract includes true, false, unknown,
conflict, out of scope and error, together with diagnostics and provenance.
The tags are not interchangeable labels for a single confidence score. The
ensemble discrepancy taxonomy must preserve them so that a failed parser is
not processed as a legal disagreement and an unresolved definition is not
processed as a transient network failure.

The precedence is part of the specification. Reordering checks can change
both a result and its explanation. The g19 example deliberately returns
conflict even though an actor scope fact is false. A Java implementation that
short-circuits on scope before consulting the required conflict closure would
disagree with the current target. That disagreement is wrong relative to the
declared execution semantics, independently of whether the bank later chooses
to revise the policy.

\section{A readable specification with a fixed translation}
\label{sec:proof-cnl}
Controlled natural language keeps the familiarity of English while restricting
how sentences are formed. The restriction matters because the reader should
be able to determine the formal structure from the sentence without asking a
model to interpret it again. Logical English provides one developed example of
this approach through explicit templates, variable binding and restrictions on
ordinary linguistic ambiguity \citep[language and law examples]{logicalenglish2023}.

For our narrow gift example, a specification can state: ``For an assessed
component in the established distributor/fund scope, the selected gift
restriction is triggered if the component is a gift, at least one of the two
specified promotional links holds, and the component does not meet the defined
discount exception.'' The specification must then define each term and its
evidentiary state. This sentence is reviewer-facing exposition. A machine
translation needs an accompanying restricted structure in which the three
conjuncts, disjunction and negation are explicit.

An appropriate authoring interface might present one condition per line with
typed connectives and named definitions. The reviewer sees the English label
and can inspect the exact expression beneath it. Free text remains available
for reasoning and caveats, but it does not secretly alter the executable
condition. A change to a definition or connective creates a new specification
revision. Editing an explanatory paragraph must not bypass the formal diff.

The controlled specification also records the legal character of the source.
A statement framed as an obligation, a prohibition, a permission, a discretion
or an expectation should retain that category. Translating all of them into a
Boolean called \texttt{allowed} would destroy distinctions before verification
begins. The current RuleIR chiefly computes conditions. A reviewed operational
policy determines how those conditions affect a bank action. Richer normative
and temporal languages are introduced in Chapter~\ref{ch:languages}.

A fixed translation from restricted syntax to RuleIR is a valuable target for
proof. Every permitted construct has a known translation rule; unsupported
sentences are rejected or left as unresolved authoring work. If a reviewer
writes an unrestricted sentence outside that syntax, a model can propose a
restricted rendering, but that proposal returns to the interpretation process.
The system cannot claim deterministic controlled-language compilation merely
because an LLM usually produces valid JSON.

Roundtrip verification supplies an additional check. Translate a formal
candidate into readable language, formalize the rendering again, and compare
the two formal candidates. A mismatch can reveal drift. Agreement can show
stability across the tested transformations. The same omitted source condition
can survive both translations, so the roundtrip result remains one component
of the evidence, as acknowledged by \citet[method and limitations]{roundtrip2026}.
It cannot replace the independent source inventory established earlier.

\section{How a preservation argument is constructed}
\label{sec:proof-induction}
The phrase ``prove the compiler correct'' becomes more useful when the reader
can see the structure of such a proof. Consider only the finite Boolean
expression grammar in equation~\eqref{eq:boolean-grammar}. Let its source
evaluator use Table~\ref{tab:kleene}. Suppose a translation maps literals,
fact references and each connective to corresponding target operations.
Assume the target fact lookup returns the same normalized tagged value for
every declared fact and that its primitive operations implement the same
tables.

\begin{proposition}[Preservation for the Boolean expression fragment]
Under those assumptions, translating any well-formed finite expression in
equation~\eqref{eq:boolean-grammar} preserves its three-valued result for every
normalized fact assignment.
\label{prop:boolean-preservation}
\end{proposition}
\begin{proof}
Use induction on the construction of the expression. A true or false literal
has the same value by the literal translation. A fact reference has the same
value by the lookup assumption. For a negation, the induction hypothesis gives
the same child value in both evaluators, and the target negation table maps
that value to the same result. For a conjunction or disjunction, the induction
hypotheses give equality for every child. Applying the identical connective
table to the identical sequence of values gives the identical parent value.
These cases cover every constructor in the grammar. The finite syntax tree
ensures the induction reaches a base case on every path.
\end{proof}

This small proof is exact, but its scope is intentionally visible. It assumes
correct target primitives and lookup. It does not establish the current Java
implementation satisfies those assumptions. It does not cover scope, conflict
normalization, numeric operations, exceptions, errors or traces until separate
lemmas include them. And it says nothing about whether the expression is a
justified interpretation of paragraph 10.

The next proof obligations follow the same pattern. A normalization lemma must
show that both implementations select the same admissible evidence under the
same times and supersession relation. A scope lemma must preserve the tagged
result and the decision to skip the body. A conflict lemma must preserve the
static closure and conflict precedence. A trace lemma must connect each
evaluated or skipped node to the required source and evidence identifiers.
Only after these claims are assembled can the implementation claim be expanded
from a Boolean value to the full semantic result.

The actual generated Java class currently binds an immutable validated bundle
and invokes the Java semantic runtime. The proof boundary must follow that
architecture. Proving a small emitted wrapper correct while assuming an
unchecked runtime would leave most of the computation outside the theorem.
Conversely, a proof about a generic runtime still needs to show that the
generated class embeds the intended bundle and that the host invokes the
correct class. The theorem's trusted components should be named rather than
hidden under the word ``compiler.''

There is a useful discipline here for an implementation agent. Begin with a
small formal fragment whose semantics and generated behavior can be inspected.
Prove or exhaustively check its components, then extend the fragment with new
constructors and corresponding obligations. A failed lemma is information
about the specification or implementation. It is not a reason to weaken the
claimed target without explaining the change.

\section{Counterexamples and the domain they live in}
\label{sec:proof-counterexamples}
To compare two candidate interpretations, the system first needs a common
question and compatible inputs. Let $F_1$ and $F_2$ compute tagged results on
the same declared input domain. A solver can search for an assignment satisfying
\begin{equation}
 D(x)\land\left(\pi(\eval(F_1,x))\ne
                         \pi(\eval(F_2,x))\right).
 \label{eq:difference-query}
\end{equation}
Here $D$ states the permitted facts and relationships. In the product-type
example, a witness sets gift true, specific-product link false, type link true
and discount false, with scope established. The assignment has different
consequences under the complete and omitted-route rules.

The first solver check should establish that $D$ is satisfiable. If the domain
simultaneously requires a fact to be true and false, the difference query is
unsatisfiable for an unhelpful reason: there are no admissible inputs. The
correct response is \texttt{INCONSISTENT\_DOMAIN}, not equivalence. This
precheck is a concrete safeguard against a formally successful but meaningless
verification result.

If the domain is satisfiable and the difference query is satisfiable, the
assignment is a counterexample to equivalence in that domain. Before displaying
it to a reviewer, replay it through the ordinary interpreters. That replay
checks that the solver encoding and the runtime agree on what the witness
means. If replay disagrees with the solver prediction, the comparison harness
has failed and the witness cannot be used to adjudicate the candidates.

If the difference query is unsatisfiable, the supported conclusion is absence
of a counterexample under the encoded semantics and domain. The strength of
that result depends on the solver procedure and any checked certificate. A
machine-checked certificate provides a different assurance level from trusting
an unchecked solver response. A timeout, unsupported construct or unknown
solver result remains unresolved; it must never be converted to successful
equivalence because no witness was printed.

Domain assumptions can hide exactly the error we hoped to find. If a comparison
asserts that every product-type promotion also names a specific product, the
omitted-route mutation may become indistinguishable from the proposed rule.
That assumption would need a source or business justification. It cannot be
introduced merely because it makes the proof succeed. The report should show
which constraints exclude a witness and where those constraints came from.

Minimal witnesses are useful for human review, but ``minimal'' also requires a
definition. One option minimizes the number of facts changed from a familiar
baseline. Another minimizes monetary amounts within meaningful bounds. Neither
guarantees that the resulting scenario is legally or operationally feasible.
Source and domain validation remain necessary. The smallest formal assignment
may be a poor real-world question if it combines classifications that a reviewer
would never jointly accept.

\section{Proving a workflow property}
\label{sec:proof-workflow}
Some of the strongest useful guarantees concern the review process rather than
the legal reading. For example, the proposed ensemble should never mark an
assessment releasable while a material discrepancy remains unresolved. This is
a property of explicit records and transitions. It is amenable to a state-machine
argument once the terms ``material,'' ``unresolved'' and ``releasable'' have
precise stored meanings.

Let $Q$ be an assessment state, $M(Q)$ its set of unresolved material issues,
and $R(Q)$ the predicate that a release is authorized. A required invariant is
\begin{equation}
 R(Q)\ \Longrightarrow\ M(Q)=\varnothing.
 \label{eq:no-open-release}
\end{equation}
The invariant does not say that an empty issue set proves the interpretation
correct. It says that known unresolved material issues cannot be ignored by
the release transition. A faulty detector may fail to create an issue, which
is why independent coverage and detector evaluation remain necessary.

To establish the invariant, the initial state has no release. Every transition
that creates, reopens or materially changes an issue must invalidate dependent
release eligibility. The only transition that authorizes release must check
the issue set and all other required conditions in the same transaction that
records the authorization. If a concurrent worker can add an issue after the
check but before the release without a version check, the implementation does
not preserve the invariant. Database transaction boundaries are therefore part
of the proof assumptions, not an incidental engineering detail.

A second property concerns termination of automatic processing. Suppose a run
has a finite number of reserved actions, and every dispatched action consumes
one reservation even if the worker fails. If all automatic paths either
consume a reservation or enter a terminal state, then no run can dispatch more
actions than its budget. This argument requires crash and retry accounting;
counting only successful responses would permit unlimited failed calls.
Chapter~\ref{ch:ensemble} gives the exact controller and its termination
measure.

These guarantees directly support the user's requirement. They can establish
that the system reports remaining uncertainty and respects its retry limit.
They cannot establish that the detectors discovered every possible uncertainty.
The two responsibilities should be tested separately: detection quality against
independent cases, and workflow integrity against state-transition and failure
injection tests.

\section{Temporal duties require histories}
\label{sec:proof-temporal}
A Boolean assessment of an offer is not enough for every circular requirement.
Some duties become active after an event and must be performed before a
deadline. Others must remain satisfied throughout an interval. Consent can be
granted, withdrawn and granted again. A formal specification for these cases
needs histories and states, not just a snapshot of current fields.

Consider a hypothetical duty activated at time $a$ with an exclusive deadline
$d$. Under a declared non-preemptive achievement interpretation, a matching
performance at time $t$ satisfies the timing condition exactly when
\begin{equation}
 a\leq t<d.
 \label{eq:performance-window}
\end{equation}
The inequality states an engineering profile; a source may require a different
boundary or permit performance before activation. Those choices must be
reviewed and encoded explicitly. The event must also match the duty's actor,
action and subject. A timely action by the wrong party is not satisfaction.

Failure to find a performance record is not by itself evidence that the duty
was breached. The observation stream may be incomplete. The current event
runtime therefore requires a completeness assertion through the deadline before
inferring breach from absence. That assertion has provenance and a recording
time. A process clock cannot create it. This mirrors the source-coverage
problem: an empty retrieved set does not prove that the relevant object does
not exist.

Late performance and breach are also distinct. A duty can be performed after
its deadline without erasing the fact that timely performance was absent under
complete evidence. A later correction may establish that an earlier performance
was missed; that creates a new assessment with a different knowledge cutoff.
The original assessment remains available. Normative logic and contract-state
languages help express these distinctions, but their semantics must be selected
to match the intended duty \citep{violations2006,normative2016,stipula2021}.

The current profile supports a limited event model. It does not automatically
implement every waiver, sanction, compensatory duty or maintenance obligation
that a legal document could contain. Unsupported temporal constructs should
remain explicit implementation gaps. A translator that invents a deadline or
converts a continuing duty into a single check would create a new rule while
appearing to simplify the software.

\section{An assurance claim is a conjunction of different evidence}
\label{sec:proof-assurance}
We can now state the overall structure without pretending that one theorem
settles the entire process. Source retention establishes which material was
reviewed. Interpretation review establishes an attributable decision about that
material within a declared context. Formal checks establish consequences and
preservation properties for specified languages and domains. Integration checks
establish how a result reaches a bank action. Each supports a distinct premise
in the justification for using a control.

This structure prevents two opposite mistakes. The first is to dismiss formal
verification because it cannot prove the meaning of unrestricted English.
Preserving a reviewed specification across implementations is valuable, and
workflow invariants can rule out important classes of operational failure. The
second is to present formal verification as if it eliminated the interpretation
and evidence premises. A perfectly verified formula still needs the right
subject, source version and fact classifications.

For the bank's reviewers, the practical outcome is a set of answerable
questions. Which parts of the interpretation are accepted and by whom? Which
formal properties have actually been established? What lies outside the
supported fragment? Which data mappings remain unresolved? What would reopen
the conclusion? Those questions can be supported by exact records rather than
an undifferentiated confidence statement.

The next chapter examines existing languages and methods in that light. Each
is introduced by the problem it solves: readable authoring, exceptions,
alternative interpretations, events, or symbolic checking. Their combination
must preserve those distinctions instead of assuming that choosing one formal
language supplies the whole assurance argument.
